\chapter{Evaluation}

\subsection{Inferenzstatistische Prüfung der ersten Modulationsannahme (M1)}

Die erste Forschungshypothese $H_{1, M1}$ postulierte, dass die Anwendung des PSI-basierten Modells ($B_{PSI}$) im Kontext einer chronischen Willenshemmung zu einer signifikant stärkeren Aktivierung der Intuitiven Verhaltenssteuerung ($IVS$) führt als die konventionelle Kontrollbedingung ($B_K$). 

Der einseitige Wilcoxon-Vorzeichen-Rang-Test für verbundene Stichproben bestätigt diese Annahme. Die Experimentalgruppe erzielte eine mittlere positive Verschiebung des $IVS$-Scores von $M = +24{,}05\%$, während die Kontrollgruppe lediglich einen Zuwachs von $M = +11{,}78\%$ aufwies. Die inferenzstatistische Analyse weist eine hochsignifikante Abweichung von der Nullhypothese auf:
\begin{math}
    W = 4753{,}0; \ z = +7{,}661; \ p < 0{,}00001
\end{math}.

Die standardisierte Effektstärke nach Cohen beträgt $r = 0{,}766$. Gemäß den Schwellenwerten von Cohen (1988) indiziert dies einen starken experimentellen Steuerungseffekt. Die Nullhypothese $H_{0, M1}$ kann somit vollständig verworfen und die Alternativhypothese $H_{1, M1}$ als empirisch bestätigt angesehen werden. Die mathematische Konditionierung des Suchraums über die berechnete Diskrepanzvariable $\tau_{m1}$ führt zu einer hochgradig robusten und praxisrelevanten Handlungsbahnung.

\subsection{Inferenzstatistische Prüfung der zweiten Modulationsannahme (M2)}

Die zweite Forschungshypothese $H_{1, M2}$ untersuchte den Transformationseffekt bei vorliegender stressinduzierter Hemmung der Selbstentwicklung. Es wurde angenommen, dass die Bedingung $B_{PSI}$ eine signifikant stärkere Aktivierung des erfahrungsbasierten Extensionsgedächtnisses ($EG$) über Mechanismen der angeleiteten Selbstberuhigung bewirkt als die Bedingung $B_K$.

Deskriptiv verzeichneten beide Gruppen eine deutliche Erhöhung des holistischen Systemanteils. Unter der PSI-basierten Dialogsteuerung stieg der relative $EG$-Anteil im Mittel um $M = +26{,}51\%$, während unter der konventionellen Bedingung ein Zuwachs von $M = +21{,}72\%$ gemessen wurde. Der mathematische Paarvergleich bestätigt die statistische Überlegenheit des Interventionsmodells:
\begin{math}
    W = 3053{,}0; \ z = +1{,}815; \ p = 0{,}01309
\end{math}.

Da der p-Wert unter dem kritischen Schwellenwert von $\alpha = 0{,}05$ liegt, wird die Nullhypothese $H_{0, M2}$ ebenfalls verworfen und die Alternativhypothese $H_{1, M2}$ angenommen. Die standardisierte Effektstärke beträgt hierbei jedoch lediglich $r = 0{,}182$, was einem schwachen experimentellen Steuerungseffekt entspricht. 

Obwohl das System $B_{PSI}$ die theoriegeleitete Selbstberuhigung und die Integration in das Selbstsystem im Vergleich zur Kontrollgruppe statistisch nachweisbar verbessert, fällt die relationale Differenzierung auf dieser Achse experimentell deutlich geringer aus als auf der Handlungsbahnungsachse ($M1$). Dies deutet darauf hin, dass konventionelle, inhärent empathische traumasensible Kommunikationsmuster (Bedingung $B_K$) bereits eine starke basale Entlastung und damit eine implizite Herabregulierung des negativen Affekts bewirken, wodurch sich der zusätzliche translationale Effekt des mathematischen Regelwerks relativiert.


\section{Bewertung des Modells}